\section{Introduction}
Large Language Models (LLMs) have demonstrated remarkable capabilities across a wide range of natural language processing tasks \citep{llm-survey}, including Qwen3 \citep{qwen3}, Kimi K2.5 \citep{k2.5}, and DeepSeek-R1 \citep{deepseek-r1}. To align these models with human intent and specific task requirements, Reinforcement Learning (RL) has become a standard paradigm \citep{rl-survey, sft-rl}. Recently, Group Relative Policy Optimization (GRPO) \citep{grpo} and its variants \citep{dapo, gspo, vcrl} have emerged as highly efficient alternatives to Proximal Policy Optimization (PPO) \citep{ppo} for LLMs. By eliminating the need for a separate value model and relying instead on relative advantage estimation within a sampled group of rollouts, GRPO significantly reduces memory overhead and simplifies the training pipeline \citep{liu2025part}.

However, deploying LLMs in real-world scenarios rarely involves optimizing a single, isolated metric. Practical applications dictate multi-objective requirements: a model must not only provide accurate answers but also adhere to length constraints \citep{efficient_survey1, efficient_survey2}, minimize bug rates in code generation \citep{DBLP:journals/ese/TambonDNKDA25, gao2025survey}, maintain a low hallucination rate \citep{hallucination_survey1, hallucination_survey2}, and keep correct tool-calling format in tool-use \citep{search_r1, airrag}. Adapting GRPO to this multi-reward setting is non-trivial. The standard practice involves scalarization-either linearly combining the raw rewards (Reward Combination) or independently normalizing the rewards and then combining their respective advantages (Advantage Combination).

Despite their widespread use, both methods suffer from significant theoretical and practical drawbacks. As we demonstrate in this work, the Reward Combination method frequently generates advantages with excessively large squared magnitudes than the Advantage Combination method, which translates to erratic policy gradients and training instability. Conversely, while the Advantage Combination method normalizes these magnitudes, it relies on static hyperparameters and completely isolates the objectives during normalization. This naive decoupling fails to capture the intricate correlations—whether synergistic or antagonistic—between different objectives during a single rollout, often leading to suboptimal trade-offs.

To address these fundamental limitations, we propose \textbf{D}ynamic \textbf{V}ariance-adaptive \textbf{A}dvantage \textbf{O}ptimization (\textbf{DVAO}). DVAO elegantly bridges the gap between stability and objective synergy by dynamically adjusting the combination weights based on the empirical reward variance of each objective within the rollout group. This completely data-driven method up-weights objectives with higher variance—indicating a stronger learning signal—while suppressing noisy, low-variance objectives. Crucially, we mathematically prove that DVAO not only bounds the advantage magnitude for stable training but also introduces a self-adaptive cross-objective regularization mechanism. In DVAO, the gradient contribution of a single objective is modulated by the overall multi-objective performance of that specific rollout, ensuring a holistic optimization trajectory.

% The main contributions of our work are summarized as follows:
% \begin{itemize}
% \item \textbf{Theoretical Analysis of Multi-Reward GRPO:} We provide a rigorous mathematical analysis exposing the critical flaws of existing scalarization methods in GRPO, demonstrating how Reward Combination leads to magnitude explosion and how Advantage Combination ignores cross-objective interactions.
% \item \textbf{Novel Optimization Method:} We introduce DVAO, a fully dynamic, hyperparameter-free weighting scheme for multi-objective GRPO. We prove theoretically that DVAO maintains bounded advantage magnitudes while introducing an implicit cross-objective regularization term that promotes synergistic learning.
% \item \textbf{Empirical Superiority [Placeholder]:} Through extensive experiments on [Insert Task/Dataset names here, e.g., code generation and safety alignment], we demonstrate that DVAO significantly outperforms baseline methods. It achieves a superior multi-objective Pareto frontier, accelerates convergence, and maintains robust training stability even when scaling to [Insert number] conflicting reward functions.
% \end{itemize}

In summary, we theoretically expose the fundamental flaws of existing scalarization methods in multi-reward GRPO—namely magnitude explosion and objective isolation—and propose Dynamic Variance-adaptive Advantage Optimization to address these limitations. DVAO is a fully dynamic, hyperparameter-free weighting scheme that we mathematically prove maintains bounded advantage magnitudes while introducing an implicit cross-objective regularization mechanism to promote synergistic learning. Extensive empirical evaluations on mathematical reasoning and tool-use benchmarks demonstrate that DVAO significantly outperforms baseline methods, accelerating convergence and consistently achieving a superior multi-objective Pareto frontier without sacrificing robust training stability.